\section{Objetivos}

\subsection{Objetivo general}

Diseñar, transformar, implementar y validar una base de datos relacional en Oracle Database 19c para la gestión de procesos odontológicos, tomando como referencia documental la tesis del sistema web LALYSDENT desarrollado con Laravel y MySQL.

\subsection{Objetivos específicos}

Identificar las entidades, atributos, relaciones, cardinalidades y restricciones del dominio odontológico documentado en la tesis base.

Elaborar el modelo conceptual en notación Chen usando diagrams.net.

Transformar el modelo conceptual al modelo relacional, definiendo tablas, claves primarias, claves foráneas, restricciones de dominio y relaciones opcionales u obligatorias.

Construir el modelo lógico en notación Crow's Foot mediante ER/Studio 2003.

Implementar el diseño físico en Oracle Database 19c, usando SQL y PL/SQL.

Crear un diccionario de datos completo y contrastarlo con el catálogo técnico de Oracle.

Comparar el enfoque MySQL/Laravel descrito en la tesis base con la transformación realizada hacia Oracle 19c.

Validar la base de datos mediante datos de prueba, consultas de control, triggers y procedimiento almacenado.

\section{Herramientas utilizadas}

\begin{table}[H]
\caption{Herramientas utilizadas en el proyecto}
\centering
\small
\begin{tabularx}{\textwidth}{L{0.28\textwidth}Y}
\toprule
\textbf{Herramienta} & \textbf{Uso en el proyecto} \\
\midrule
diagrams.net / Draw.io & Elaboración del modelo conceptual en notación Chen. \\
ER/Studio 2003 & Elaboración del modelo lógico relacional en notación Crow's Foot. \\
Oracle Database 19c & Motor de base de datos usado para la implementación física. \\
Docker & Ejecución de Oracle 19c en contenedor. \\
DBeaver & Ejecución de scripts SQL/PLSQL y validación de objetos. \\
VSCodium & Edición del proyecto, scripts y documento final. \\
Konsole con zsh & Ejecución de comandos, Git, compilación LaTeX y organización del repositorio. \\
LuaLaTeX + Biber & Compilación del informe final con formato APA 7. \\
Git & Control de versiones del proyecto. \\
\bottomrule
\end{tabularx}
\end{table}

\section{Análisis de requisitos}

\subsection{Materiales del proyecto y fuentes de apoyo}

Para contextualizar el dominio LALYSDENT se consideraron dos materiales principales. El primero fue el sitio web de la tesis base, disponible en \url{https://calvocobos.github.io/}, usado como apoyo documental para comprender los procesos de registro, atención, inventario y finanzas del consultorio odontológico. El segundo fue el repositorio Laravel desarrollado por el equipo del proyecto, disponible en \url{https://github.com/AlanDevPro/sistema-odontologico}, usado para ejecutar y evidenciar la aplicación web propia del sistema \parencite{calvocobos2025web,sistemaodontologicogithub}.

La aplicación Laravel forma parte del trabajo del equipo, pero no reemplaza el objetivo principal del examen. El producto central de este informe sigue siendo el diseño, transformación, implementación y validación de la base de datos Oracle 19c en el schema \texttt{FARMACIAS\_BD3}. Para ejecutar la aplicación web sin alterar dicho diseño, se utilizó un schema separado llamado \texttt{LALYSDENT\_APP}.

\subsection{Alcance funcional tomado de la tesis base}

La tesis base desarrolla un sistema de información web para administrar procesos de registro, atención, inventario y finanzas del consultorio odontológico LALYSDENT. Dicho sistema incluye módulos de registro de pacientes, gestión de citas, historias clínicas con odontograma, administración de tratamientos, control de suministros e inventario odontológico, pagos de pacientes y reportes consolidados \parencite{tesisbase,calvocobos2025web}.

En este examen final no se documenta el frontend ni la arquitectura completa de la aplicación web. El alcance se limita a la base de datos: diseño conceptual, transformación relacional, modelo lógico, implementación física, diccionario de datos, validación y comparación técnica.

\subsection{Entidades principales}

A partir del dominio se identificaron las siguientes entidades: asistente, folder, doctor, tratamiento, proveedor, suministro, paciente, cita, odontograma, detalle de odontograma, plan de pago, pago, compra, detalle de compra e inventario.

\subsection{Procesos cubiertos}

La base de datos cubre cuatro procesos principales:

\begin{itemize}
  \item Registro y citas: pacientes, carpetas clínicas, doctores, asistentes y citas.
  \item Atención odontológica: odontogramas, detalle por pieza dental y tratamientos.
  \item Finanzas y pagos: planes de pago y pagos registrados.
  \item Compras e inventario: proveedores, suministros, compras, detalle de compra y stock.
\end{itemize}

\subsection{Supuestos de transformación}

No se dispuso de un DDL MySQL original completo de la tesis. Por ello, la comparación y transformación se realizaron a partir de la evidencia documental de la tesis y de los módulos funcionales descritos. La transformación no afirma una equivalencia sentencia por sentencia con un script MySQL original, sino una conversión técnica documentada del modelo de datos hacia Oracle 19c.
